\scanpage{278}


et il s'agit de voir que l'on \unsure{ne peut|en peut} avoir

\[
u=e+f,\quad u'=e-f
\]

sauf \deleted{dans le cas précisé} si $2\,1_\Lambda=0$, (où on a $\inserted{+1_\Lambda=-1_\Lambda\ \text{donc}}u=u'$, et on gagne), \inserted{ou dans le cas précisé dans l'énoncé ci-dessus.} Supposons donc

\[
2.1_\Lambda\ne0
\]

Notons que \deleted{puisque $u\equiv u'$}\inserted{comme $u_i^\circ,{u'}_i^\circ$} sont colinéaires dans tout corps résiduel $k_i$, cel implique déjà que $2\,1_{k_i}=0$ pour tout $i$, i.e.

\[
2.1_\Lambda\in\unsure{\mathfrak r|\mathfrak R|k}\,\Lambda.
\]

La condition $(*)$ sur $x\in M^*$ s'écrit, écrivant

\[
x=\alpha e+\beta f\qquad\alpha,\beta\in\Lambda
\]

que

\[
\alpha+\beta\in\Lambda^*\ \Longrightarrow\ \deleted{\unsure{2\alpha\ \text{ou}|2\alpha\ \text{ou }2\beta}}\alpha+\beta\in\pm(\alpha-\beta)\quad\text{i.e. }2\alpha\ \text{ou }2\beta\ \text{nul.}
\]

Posant $\lambda=(\alpha+\beta)$ et multipliant $\alpha,\beta$ par $\lambda^{-1}$, on est ramené à voir que

\[
\begin{aligned}
&\alpha,\beta\in\Lambda,\ \alpha+\beta=1\Longrightarrow 2\alpha\ \text{ou }2\beta\ \text{nul }\deleted{\unsure{\text{ce}}},\ \text{i.e.}\\
&\forall\,\alpha\in\Lambda,\ \text{on a }2\alpha\ \text{ou }2(1-\alpha)=2-2\alpha\ \text{nul, i.e.}\\
&\deleted{\unsure{(**)}}\qquad 2\Lambda\subset\{0,2\},\quad\text{donc (puisque }2.1_\Lambda\ne0\text{)}\\
&\deleted{\unsure{(***)}}\qquad \boxed{2\Lambda=\{0,2\}\ne\{0\}},\quad\text{(qui implique }4_\Lambda=0\ \ldots\text{)}
\end{aligned}
\]

\uline{Remarque}

On voit que dans ce cas, les \deleted{couples}\inserted{paires} $\{u,u'\}$ qui mettent en défaut la conclusion escomptée $\{u'\in\pm u\}$ forment une orbite sous $\operatorname{Gl}_\Lambda(M)$, On peut aussi écrire, plus précisément

\[
\begin{aligned}
u'&=u-2f=u+2f\qquad(\text{NB }4_\Lambda=0),\\
(\text{d'où}\quad -u'&=-u-2f=-u+2f)
\end{aligned}
\]

où $2f$ ne dépend que de $f$ modulo l'idéal \inserted{$\mathfrak m_{i_0}$} annulateur de 2, qui est un idéal maximal à corps résiduel \deleted{\unsure{$\mathbb F_2$}}

\[
k_{i_0}\simeq\mathbb F_2
\]

on a \deleted{\unsure{donc}} une suite exacte

\[
0\xrightarrow{\unsure{1|{}}}\underset{\substack{\simeq 2M\\\text{vectoriel de}\\\text{rang 2 sur }k_{i_0}\simeq\mathbb F_2}}{\underbrace{M\otimes_\Lambda k_{i_0}}}\longrightarrow M\longrightarrow M/2M\longrightarrow0
\]
% archive sheet number: 277
